\section{Preliminaries}
\label{sec:prelim}
% -----------------------------------------------------------------------

This section lays the groundwork for the rest of the paper and fixes the
optimization framework, model-structural conventions, and mathematical notation
used in the analysis that follows.  Section~\ref{sec:prelim-formulation} first
formalizes large-language-model training as a stochastic optimization problem
and clarifies its characteristic parameter topology.  Building on this,
Section~\ref{sec:prelim-background} reviews the core mechanisms of classical
optimization methods, ranging from SGD and momentum to adaptive moment
estimation.  It also introduces the preconditioning perspective behind the
Newton method, the natural gradient, and Fisher or Hessian curvature, so that the
later sections have the background they need to develop the unified theoretical
framework.

\subsection{Problem Formulation}
\label{sec:prelim-formulation}

\paragraph{Training objective and stochastic gradients.}
Let $\mathcal{D}=\{(\mathbf{x}_i,\mathbf{y}_i)\}_{i=1}^N$ denote a training
corpus and let $\theta\in\mathbb{R}^d$ collect all trainable parameters of a
language model.  Pretraining and supervised fine-tuning both reduce to empirical
risk minimization,
\begin{equation}
  \min_{\theta\in\mathbb{R}^d}\;\mathcal{L}(\theta)
  = \frac{1}{N}\sum_{i=1}^{N}\ell(\theta;\,\mathbf{x}_i,\mathbf{y}_i),
  \label{eq:objective}
\end{equation}
where $\ell$ is usually a token-level cross-entropy or an instruction-tuning
loss derived from it. At step $t$, the optimizer receives a mini-batch
gradient,
\begin{equation}
  g_t = \frac{1}{|\mathcal{B}_t|}\sum_{i\in\mathcal{B}_t}
        \nabla_\theta\ell(\theta_t;\,\mathbf{x}_i,\mathbf{y}_i),
  \label{eq:mbgrad}
\end{equation}
where $\mathcal{B}_t\subset\{1,\ldots,N\}$.  When a stochastic-oracle model is
needed, we write $\mathbb{E}[g_t\mid\theta_t]=\nabla\mathcal{L}(\theta_t)$ and
$\mathbb{E}[\|g_t-\nabla\mathcal{L}(\theta_t)\|^2]\le\sigma^2$.  In practice,
$g_t$ reflects data sampling, sequence packing, distributed reduction, mixed
precision, gradient accumulation, clipping, and loss scaling.  These system
details are not separate from optimizer behavior.  In fact, they are what
determine the noise level, numerical range, and memory traffic seen by the
update rule.

\paragraph{Transformer parameter topology.}
Transformer-style LLMs~\cite{vaswani2017attention} contain heterogeneous
parameter tensors.  For optimizer design, the most important distinction is not
the layer name but the tensor geometry.  Two-dimensional matrices
$W\in\mathbb{R}^{m\times n}$, such as attention projections
$W_Q,W_K,W_V,W_O$, feed-forward projections, and embedding tables, admit
row-column operations including singular-value decompositions, Kronecker
factorization, matrix orthogonalization, and low-rank projection.  Vector-like
parameters, such as biases and normalization gains or shifts, do not have the
same matrix geometry and are usually routed to element-wise rules or excluded
from weight decay.  Some methods introduce finer scopes, e.g., attention-head
blocks, per-layer groups, or module-specific rules.  This parameter topology is
the basis for the scoping-and-routing stage of the universal meta-pipeline in
Section~\ref{sec:meta-pipeline}.

\paragraph{Memory budget under mixed-precision training.}
Under BF16 forward and backward computation with FP32 master weights, the
training memory can be estimated as
\begin{equation}
  M_{\mathrm{train}}
  \approx 4d + 2d + 2d + S_{\mathrm{opt}}d + M_{\mathrm{act}}.
  \label{eq:memory-budget}
\end{equation}
The first three terms correspond to FP32 master parameters, BF16 model copies,
and BF16 gradients, $S_{\mathrm{opt}}$ is the optimizer-state footprint in bytes
per parameter, and $M_{\mathrm{act}}$ denotes activation and temporary-buffer
memory.  Adam stores two FP32 state tensors $m_t$ and $v_t$, so each parameter
carries an extra 8 bytes.  A 7B-parameter model therefore needs about 56\,GB for
Adam states alone, before activations, temporary buffers, and distributed-training
overheads.  State-compressed methods in Section~\ref{sec:t4} are therefore not
merely engineering variants.  Under a fixed hardware budget, reducing
$S_{\mathrm{opt}}$ changes the feasible model size, batch size, and tuning regime.

\subsection{Classical Optimization Background}
\label{sec:prelim-background}

The modern LLM optimizer literature can be read as a sequence of modifications
to a small set of classical design choices.  These choices are whether to
smooth gradients, whether to maintain a curvature proxy, how to scale the
direction, and how to regularize the final write.  The following review fixes the baseline equations
used by later sections.


\paragraph{Newton method and the perspective of preconditioning.}
A natural starting point for our subsequent preconditioning analysis is the Newton method. Near the current iterate, a second-order Taylor expansion yields:
\begin{equation}
  \mathcal{L}(\theta_t+\delta)
  \approx
  \mathcal{L}(\theta_t)
  + g_t^\top\delta
  + \frac{1}{2}\delta^\top H_t\delta,
  \quad
  H_t=\nabla^2\mathcal{L}(\theta_t).
  \label{eq:newton-quadratic}
\end{equation}
Minimizing this local quadratic model gives the damped Newton update
\begin{equation}
  \theta_{t+1}
  = \theta_t - \eta_t H_t^{-1}g_t,
  \label{eq:newton-update}
\end{equation}
where the scalar $\eta_t$ absorbs damping, line search, or trust-region control.  This equation gives the classical preconditioning viewpoint, since the optimizer does not only choose a gradient signal but also chooses the geometry in which that signal is scaled.  SGD can be read as the first-order simplification that replaces $H_t^{-1}$ by the identity.  Modern large-model optimizers replace the exact inverse Hessian with tractable approximations or geometry-induced direction operators, where Adam uses a diagonal second-moment or Fisher-like proxy, Shampoo and SOAP use structured matrix factors or eigenbases, and Muon uses a polar or Gram-induced matrix direction.  This connection motivates the preconditioned and matrix-structured views used throughout the rest of the survey.

\paragraph{SGD and momentum.}
Stochastic gradient descent (SGD) writes the raw mini-batch gradient into the
parameters:
\begin{equation}
  \theta_{t+1} = \theta_t - \eta_t\,g_t.
  \label{eq:sgd}
\end{equation}
Momentum replaces the instantaneous gradient with an exponential moving
average (EMA):
\begin{equation}
  m_t = \beta_1\,m_{t-1} + (1-\beta_1)\,g_t,\quad
  \theta_{t+1} = \theta_t - \eta_t\,m_t.
  \label{eq:sgdm}
\end{equation}
SGD and momentum use one global scale, or a schedule of global scales, for all
coordinates.  They therefore do not adapt to coordinate-wise gradient
magnitudes or curvature proxies.
In the notation of Eq.~\eqref{eq:newton-update}, SGD corresponds to the
identity preconditioner, $H_t=I$, so $H_t^{-1}g_t=g_t$.  Momentum keeps the
same identity geometry but replaces the instantaneous signal $g_t$ with the
smoothed signal $m_t$.

\paragraph{AdaGrad and RMSProp.}
AdaGrad~\cite{duchi2011adaptive} introduced a diagonal, history-dependent
preconditioner by accumulating squared gradients:
\begin{equation}
  G_t = \sum_{\tau=1}^{t}g_\tau^2,\quad
  \theta_{t+1}
  = \theta_t - \eta_t\,\frac{g_t}{\sqrt{G_t}+\epsilon}.
  \label{eq:adagrad}
\end{equation}
Here and below, powers and divisions are element-wise unless stated otherwise.
The monotone accumulator is useful for sparse features but can make effective
learning rates decay too aggressively in long training runs.  RMSProp
~\cite{tieleman2012rmsprop} replaces the cumulative sum with an EMA:
\begin{equation}
  v_t = \beta_2\,v_{t-1}+(1-\beta_2)\,g_t^2,\quad
  \theta_{t+1}
  = \theta_t - \eta_t\,\frac{g_t}{\sqrt{v_t}+\epsilon}.
  \label{eq:rmsprop}
\end{equation}
Equivalently, AdaGrad uses the effective diagonal metric
$H_t^{\mathrm{AdaGrad}}=\operatorname{diag}(\sqrt{G_t}+\epsilon)$, whereas
RMSProp uses $H_t^{\mathrm{RMSProp}}=\operatorname{diag}(\sqrt{v_t}+\epsilon)$.
Both should be read as cheap diagonal preconditioners, that is, approximate
proxies for Fisher or Hessian curvature.

\paragraph{Adam.}
Adam~\cite{kingma2015adam} combines momentum and RMSProp-style second moments,
then corrects the initialization bias of both EMAs:
\begin{align}
  m_t &= \beta_1\,m_{t-1}+(1-\beta_1)\,g_t,\quad v_t = \beta_2\,v_{t-1}+(1-\beta_2)\,g_t^2,
  \label{eq:adam-moments}\\
  \hat{m}_t &= m_t/(1-\beta_1^t),\quad \hat{v}_t = v_t/(1-\beta_2^t),
  \label{eq:adam-bias}\\
  \theta_{t+1} &= \theta_t - \eta_t\,\hat{m}_t/(\sqrt{\hat{v}_t}+\epsilon).
  \label{eq:adam-update}
\end{align}
The diagonal preconditioner of Adam remains the reference point for most LLM
optimizers.  Later methods may alter the moment time scale, replace the
gradient estimator, change the direction map, compress the state, or add a
wrapper around the final update.  Nevertheless, they are often compared against
this Adam-like template.
In the same preconditioning notation, Adam uses the effective signal
$\hat{m}_t$ and the diagonal metric
$H_t^{\mathrm{Adam}}=\operatorname{diag}(\sqrt{\hat{v}_t}+\epsilon)$, so the
preconditioned direction is $H_t^{-1}\hat{m}_t
=\hat{m}_t/(\sqrt{\hat{v}_t}+\epsilon)$.

\paragraph{AdamW: decoupled weight decay.}
In Adam, adding $\ell_2$ regularization to the loss causes the regularization
gradient to be divided by the adaptive denominator. AdamW
~\cite{loshchilov2019decoupled} decouples weight decay from this adaptive
scaling:
\begin{equation}
  \theta_{t+1} = \theta_t
    - \eta_t\,\frac{\hat{m}_t}{\sqrt{\hat{v}_t}+\epsilon}
    - \eta_t\,\lambda\,\theta_t .
  \label{eq:adamw}
\end{equation}
This change is small at the formula level but important as a baseline convention: throughout the survey, an ``Adam-like'' optimizer means a method whose main state evolution is inherited from AdamW unless stated otherwise.
AdamW retains the Adam diagonal preconditioner $H_t^{\mathrm{Adam}}$ for data-gradient components, moving weight decay outside the preconditioned update direction and incorporating it directly in the final parameter assignment step.

\paragraph{Natural gradient and the Fisher information matrix.}
The natural gradient~\cite{amari1998natural} motivates many preconditioned
updates.  For a parameterized distribution $p_\theta(z)$, the Fisher
information matrix is:
\begin{equation}
  F(\theta)=\mathbb{E}_{z\sim p_\theta}\!\left[
    \nabla_\theta\log p_\theta(z)\,
    \nabla_\theta\log p_\theta(z)^\top
  \right].
  \label{eq:fisher-matrix}
\end{equation}
The corresponding update
$\theta_{t+1}=\theta_t-\eta_t F(\theta_t)^{-1}\nabla\mathcal{L}(\theta_t)$
uses the geometry of the predictive distribution, whereas ordinary gradient
descent uses the Euclidean geometry of raw parameters.  Exact storage and inversion of $F$ are infeasible
for modern LLMs, so practical optimizers use diagonal, block-diagonal,
Kronecker, low-rank, or stochastic approximations.
Thus, natural gradient follows the same Newton-style template, but it
replaces $H_t$ with the Fisher metric $F(\theta_t)$.  This makes it a useful
bridge to later matrix and curvature-aware methods, which differ mainly in how
they approximate, factor, or simplify this metric.

\paragraph{Gradient covariance and curvature.}
For negative log-likelihood objectives, the outer product of score gradients
defines the Fisher matrix.  Under standard regularity conditions, the gradient
covariance, the Fisher matrix, and the Hessian are approximately equal, as
given by
\begin{equation}
  \mathbb{E}\!\left[g_t g_t^\top\mid\theta_t\right]
  \approx F(\theta_t),
  \quad
  F(\theta_t) \approx H(\theta_t)=\nabla^2\mathcal{L}(\theta_t).
  \label{eq:fisher-identity}
\end{equation}
This distinction matters for a survey.  The Adam $v_t=\mathrm{EMA}(g_t^2)$ serves as a diagonal gradient-covariance estimate.  The Shampoo factors $L_t=\mathrm{EMA}(G_tG_t^\top)$ and $R_t=\mathrm{EMA}(G_t^\top G_t)$ form a Kronecker-style curvature proxy~\cite{gupta2018shampoo}, with $H_t^{\mathrm{Shampoo}}\approx L_t^{1/4}\otimes R_t^{1/4}$. Sophia estimates a clipped diagonal Hessian using stochastic second-order information~\cite{liu2023sophia}.  Section~\ref{sec:lmo-four-axis} turns these choices into Axis~II (state and curvature estimator) and Axis~III
(geometry and precondition operator) of the LMO-driven four-axis decomposition.

\paragraph{Notation.}
Table~\ref{tab:notation} collects the mathematical symbols used throughout
the paper and closes the preliminary setup before the meta-pipeline is
introduced.

% \begin{table}[htbp]
%   \centering
%   \caption{Mathematical notation.}
%   \label{tab:notation}
%   \small
%   \setlength{\tabcolsep}{10mm}
%   \begin{tabular}{@{}ll@{}}
%     \toprule
%     \textbf{Symbol} & \textbf{Meaning} \\
%     \midrule
%     $d$ & Total number of model parameters \\
%     $\theta_t\in\mathbb{R}^d$ & Parameter vector at step $t$ \\
%     $W_t\in\mathbb{R}^{m\times n}$ & Two-dimensional weight matrix \\
%     $\mathcal{L}(\theta)$ & Training loss \\
%     $g_t$ & Mini-batch stochastic gradient (vector) \\
%     $G_t\in\mathbb{R}^{m\times n}$ & Gradient in matrix form \\
%     $m_t$ & First-order momentum (gradient EMA) \\
%     $v_t$ & Second-order moment estimate (squared-gradient EMA) \\
%     $\eta_t$ & Learning rate at step $t$ \\
%     $\beta_1,\,\beta_2$ & EMA decay coefficients \\
%     $\lambda$ & Weight-decay coefficient \\
%     $\epsilon$ & Numerical stability constant \\
%     $\odot$ & Element-wise (Hadamard) product \\
%     $\langle A,B\rangle=\mathrm{tr}(A^\top B)$ & Matrix inner product \\
%     $\|\cdot\|_p$ & Vector $\ell_p$ norm \\
%     $\|\cdot\|_F$ & Frobenius norm \\
%     $\|\cdot\|_\mathrm{op}$ & Spectral (operator) norm \\
%     $\|\cdot\|_\ast$ or $\|\cdot\|_*$ & Nuclear norm for matrices; dual norm for a generic norm \\
%     $[\cdot]^\sharp$ & Dual map (sharp operator) \\
%     $F(\theta)$ & Fisher information matrix \\
%     $H(\theta)=\nabla^2\mathcal{L}(\theta)$ & Hessian matrix \\
%     $\hat{H}_t$ & Hessian estimator (form depends on Axis~II) \\
%     $\mathrm{EMA}_\beta(x_t)$ & Exp.\ moving avg.: $\beta(\cdot)_{t-1}+(1-\beta)x_t$ \\
%     $W=U\Sigma V^\top$ & Singular value decomposition of $W$ \\
%     $Q_L,\,Q_R$ & Analysis bases, left and right (LMO-driven Axis~I) \\
%     $\Phi_\mathrm{dir}$ & Direction function (LMO-driven Axis~III) \\
%     $\alpha\in[0,1]$ & Hessian exponent (LMO-driven Axis~III) \\
%     $M_t\in\mathbb{R}^{m\times n}$ & Momentum in matrix form \\
%     $\mathrm{lmo}_{\mathcal{C}}(s)$ & Linear minimization oracle over convex set $\mathcal{C}$ \\
%     $\rho^{(i)}$ & Routing label of parameter $i$ (meta-pipeline Step~1) \\
%     $\tilde{G}_t$ & Variance-reduced gradient estimate (LMO-driven Axis~II) \\
%     $T_1$--$T_5$ & Five methodological optimizer families (Dimension~A) \\
%     $O_1$--$O_6$ & Six effect objectives (Dimension~B) \\
%     \bottomrule
%   \end{tabular}
% \end{table}

\begin{table}[!t]
  \centering
  \caption{Summary of the mathematical notations used throughout this survey, including optimization variables, matrix operators, second-order quantities, and framework-specific symbols. The table serves as a reference for the formulations presented in the following sections.}
  \label{tab:notation}
  \footnotesize
  \setlength{\tabcolsep}{2pt}
  \begin{tabularx}{\textwidth}{@{}>{\raggedright\arraybackslash}p{0.17\textwidth}>{\raggedright\arraybackslash}X|>{\raggedright\arraybackslash}p{0.17\textwidth}>{\raggedright\arraybackslash}X@{}}
    \toprule
    \textbf{Symbol} & \textbf{Meaning} &
    \textbf{Symbol} & \textbf{Meaning} \\
    \midrule

    $d$ & Total number of model parameters &
    $W_t\in\mathbb{R}^{m\times n}$ & Weight matrix \\

    $\theta_t\in\mathbb{R}^{d}$ & Parameter vector &
    $G_t\in\mathbb{R}^{m\times n}$ & Gradient matrix \\

    $\mathcal{L}(\theta)$ & Training loss &
    $M_t\in\mathbb{R}^{m\times n}$ & Matrix momentum \\

    $g_t$ & Mini-batch stochastic gradient &
    $W=U\Sigma V^\top$ & Singular value decomposition \\

    $m_t$ & First-order momentum &
    $\langle A,B\rangle=\mathrm{tr}(A^\top B)$ & Matrix inner product \\

    $v_t$ & Second-order moment estimate &
    $\odot$ & Hadamard product \\

    $\eta_t$ & Learning rate &
    $\|\cdot\|_p$ & Vector $\ell_p$ norm \\

    $\beta_1,\beta_2$ & EMA decay coefficients &
    $\|\cdot\|_F$ & Frobenius norm \\

    $\lambda$ & Weight-decay coefficient &
    $\|\cdot\|_{\mathrm{op}}$ & Spectral norm \\

    $\epsilon$ & Numerical stability constant &
    $\|\cdot\|_\ast$ or $\|\cdot\|_*$ & Nuclear or dual norm \\

    $F(\theta)$ & Fisher information matrix &
    $Q_L,Q_R$ & Axis-I coordinate bases \\

    $H(\theta)=\nabla^2\mathcal{L}(\theta)$ & Hessian matrix &
    $\Phi_t$ & Axis-III direction operator \\

    $\hat{H}_t$ & Axis-II curvature proxy &
    $\alpha\in[0,1]$ & Preconditioner exponent \\

    $\mathrm{EMA}_\beta(x_t)$ & Exponential moving average &
    $\mathrm{lmo}_{\mathcal{C}}(s)$ & Linear minimization oracle \\

    $\tilde{G}_t$ & Variance-reduced gradient estimate &
    $\rho^{(i)}$ & Meta-pipeline routing label \\

    $[\cdot]^\sharp$ & Dual map &
    $T_1$--$T_5$ & Dimension-A method families \\

    && $O_1$--$O_6$ & Dimension-B effect objectives \\

    \bottomrule
  \end{tabularx}
\end{table}
